\documentclass[11pt,a4paper]{article}
\usepackage[utf8]{inputenc}
\usepackage[T1]{fontenc}
\usepackage{amsfonts}
\usepackage[french]{babel}
\frenchbsetup{StandardLists=true} % à inclure si on utilise \usepackage[francais]{babel}
\usepackage{enumitem}
\usepackage{amssymb}
\usepackage{amsmath}
\usepackage[left=2cm,right=2cm,top=2cm,bottom=2cm]{geometry}
\usepackage{multirow}
\usepackage[dvipsnames]{xcolor}

\usepackage[T1]{fontenc}
\usepackage{fouriernc}
\usepackage{graphicx}

\usepackage{setspace}
\singlespacing

\usepackage{multicol}

\usepackage{pstricks,pst-plot}
\usepackage{tikz}
\usepackage{tkz-tab}
\usepackage{tkz-base}
\usepackage{pgfplots,mathtools}
\pgfplotsset{compat=1.12}

\usepackage{fancyhdr}
\pagestyle{fancy}

\renewcommand\headrulewidth{1pt}
\fancyhead[L]{Lycée Denis Diderot }
\fancyhead[R]{Classe 1BTS}
\setcounter{page}{1}\fancyfoot[C]{Mini Contrôle}
\fancyfoot[R]{\today}
\fancyfoot[L]{Alexis RUIZ}
\usepackage{multirow,tabularx,booktabs}
\newcommand{\cadreblanc}[1]{%
\medskip\framebox[\linewidth][c]{%
\rule{0mm}{#1mm}}\par
\medskip}

\setlength{\columnseprule}{.25mm}

\renewcommand{\arraystretch}{2}
\newcommand*\acocher{\quitvmode{\fboxsep0pt \fboxrule0.8pt \fbox{\vrule height1.8ex depth.1ex width0pt \vrule height0pt depth0pt width1.9ex }}}
\DeclareMathOperator{\e}{e}

\begin{document}
\center \LARGE{MINI CONTRÔLE - 3.3 (15 minutes)}\
\\[0,3cm]

\normalsize

\hrule\
\\[0,2cm]

\flushleft \textbf{NOM :\\
Prénom:\\[0,2cm]
}
\hrule\
\\[0,2cm]

\textbf{Exercice 1 :} Déterminer la nature des fonctions définies par leur représentation graphique. Les questions sont présentées sous forme d’un QCM, il peut y avoir 1 ou 2 réponses correctes.\\ En cas de périodicité, déterminer également la valeur de la période $T$.\\ [5mm]


\begin{multicols}{2}
\begin{tikzpicture}
	% Dimensions du repere
	\def\xmin{-3.5} \def\xmax{3.55} \def\ymin{-1.25} \def\ymax{1.2}
	% Styles des axes et de la grille
	
	\tikzstyle{axe}=[->,>=stealth']
	\tikzstyle{grille}= [step=1, very thin]
	% Grille
	\draw [grille] (\xmin,\ymin) grid (\xmax,\ymax); 
	% Annotations axes et unites
	\draw [axe, thick] (\xmin,0)--(\xmax,0) node[above left]  {$t$};
	\draw [axe, thick] (0,\ymin)--(0,\ymax) node[below right] {$y$};
	\foreach \x in {-3,-2,-1,0,1,2,3} 
		\draw (\x,1pt)--(\x,-1pt) node [below right]{$\x$};	
	\foreach \y in {-1,1} 
		\draw (1pt,\y)--(-1pt,\y) node [left]{$\y$};	
	
	\draw[samples=400][ color=ForestGreen] [domain=-3.5:3.55][line width=1.5pt] plot[red](\x,{sin(3.14*\x r)});

	

	% Clip pour que les figures ne sortent pas du cadre
	%\clip (\xmin,\ymin) rectangle (\xmax,\ymax); 
\end{tikzpicture}

\acocher~~~~ fonction paire \\
\acocher~~~~ fonction impaire \\ 
\acocher~~~~ fonction périodique ~~~~ $T= ..............$ \\ 
\acocher~~~~ Aucune réponse possible \\


\begin{tikzpicture}
	% Dimensions du repere
	\def\xmin{-3.5} \def\xmax{3.55} \def\ymin{-1.25} \def\ymax{1.2}
	% Styles des axes et de la grille
	
	\tikzstyle{axe}=[->,>=stealth']
	\tikzstyle{grille}= [step=1, very thin]
	% Grille
	\draw [grille] (\xmin,\ymin) grid (\xmax,\ymax); 
	% Annotations axes et unites
	\draw [axe, thick] (\xmin,0)--(\xmax,0) node[above left]  {$t$};
	\draw [axe, thick] (0,\ymin)--(0,\ymax) node[below right] {$y$};
	\foreach \x in {-3,-2,-1,0,1,2,3} 
		\draw (\x,1pt)--(\x,-1pt) node [below right]{$\x$};	
	\foreach \y in {-1,1} 
		\draw (1pt,\y)--(-1pt,\y) node [left]{$\y$};	
	
	\draw[samples=400][ color=black] [domain=-3.5:3.55][line width=1.5pt] plot[red](\x,{cos((1.5*pi)*(\x) r)});

	

	% Clip pour que les figures ne sortent pas du cadre
	%\clip (\xmin,\ymin) rectangle (\xmax,\ymax); 
\end{tikzpicture}


\acocher~~~~ fonction paire \\
\acocher~~~~ fonction impaire \\ 
\acocher~~~~ fonction périodique ~~~~ $T= ..............$ \\ 
\acocher~~~~ Aucune réponse possible \\


\end{multicols}



\begin{multicols}{2}
\begin{tikzpicture}
	% Dimensions du repere
	\def\xmin{-3.5} \def\xmax{3.55} \def\ymin{-1.25} \def\ymax{2.2}
	% Styles des axes et de la grille
	
	\tikzstyle{axe}=[->,>=stealth']
	\tikzstyle{grille}= [step=1, very thin]
	% Grille
	\draw [grille] (\xmin,\ymin) grid (\xmax,\ymax); 
	% Annotations axes et unites
	\draw [axe, thick] (\xmin,0)--(\xmax,0) node[above left]  {$t$};
	\draw [axe, thick] (0,\ymin)--(0,\ymax) node[below right] {$y$};
	\foreach \x in {-3,-2,-1,0,1,2,3} 
		\draw (\x,1pt)--(\x,-1pt) node [below right]{$\x$};	
	\foreach \y in {-1,1} 
		\draw (1pt,\y)--(-1pt,\y) node [left]{$\y$};	
	
	\draw[samples=400][ color=red] [domain=0:2][line width=1.5pt] plot(\x,{2-\x});
	\draw[samples=400][ color=red] [domain=-2:0][line width=1.5pt] plot(\x,{\x+2});
	\draw[samples=400][ color=red] [domain=2:3.55][line width=1.5pt] plot(\x,{\x-2});
	\draw[samples=400][ color=red] [domain=-3.55:-2][line width=1.5pt] plot(\x,{-\x-2});

	% Clip pour que les figures ne sortent pas du cadre
	%\clip (\xmin,\ymin) rectangle (\xmax,\ymax); 
\end{tikzpicture}

\acocher~~~~ fonction paire \\
\acocher~~~~ fonction impaire \\ 
\acocher~~~~ fonction périodique ~~~~ $T= ..............$ \\ 
\acocher~~~~ Aucune réponse possible \\


\begin{tikzpicture}
	% Dimensions du repere
	\def\xmin{-3.5} \def\xmax{3.55} \def\ymin{-1.55} \def\ymax{1.55}
	% Styles des axes et de la grille
	
	\tikzstyle{axe}=[->,>=stealth']
	\tikzstyle{grille}= [step=1, very thin]
	% Grille
	\draw [grille] (\xmin,\ymin) grid (\xmax,\ymax); 
	% Annotations axes et unites
	\draw [axe, thick] (\xmin,0)--(\xmax,0) node[above left]  {$t$};
	\draw [axe, thick] (0,\ymin)--(0,\ymax) node[below right] {$y$};
	\foreach \x in {-3,-2,-1,0,1,2,3} 
		\draw (\x,1pt)--(\x,-1pt) node [below right]{$\x$};	
	\foreach \y in {-1,1} 
		\draw (1pt,\y)--(-1pt,\y) node [left]{$\y$};	
	
	\draw[samples=400][ color=blue] [domain=-3.5:-3][line width=1.5pt] plot[red](\x,{0});
    \draw[samples=400][ color=blue] [domain=-3:-2][line width=1.5pt] plot[red](\x,{-1});
    \draw[samples=400][ color=blue] [domain=-2:-1][line width=1.5pt] plot[red](\x,{1});
	\draw[samples=400][ color=blue] [domain=-1:0][line width=1.5pt] plot[red](\x,{0});
	\draw[samples=400][ color=blue] [domain=0:1][line width=1.5pt] plot[red](\x,{-1});
    \draw[samples=400][ color=blue] [domain=1:2][line width=1.5pt] plot[red](\x,{1});
    \draw[samples=400][ color=blue] [domain=2:3][line width=1.5pt] plot[red](\x,{0});
    \draw[samples=400][ color=blue] [domain=3:3.5][line width=1.5pt] plot[red](\x,{-1});
    
	% Clip pour que les figures ne sortent pas du cadre
	%\clip (\xmin,\ymin) rectangle (\xmax,\ymax); 
\end{tikzpicture}


\acocher~~~~ fonction paire \\
\acocher~~~~ fonction impaire \\ 
\acocher~~~~ fonction périodique ~~~~ $T= ..............$ \\ 
\acocher~~~~ Aucune réponse possible \\


\end{multicols}

\begin{multicols}{2}
\begin{tikzpicture}
	% Dimensions du repere
	\def\xmin{-3.5} \def\xmax{3.55} \def\ymin{-1.25} \def\ymax{1.2}
	% Styles des axes et de la grille
	
	\tikzstyle{axe}=[->,>=stealth']
	\tikzstyle{grille}= [step=1, very thin]
	% Grille
	\draw [grille] (\xmin,\ymin) grid (\xmax,\ymax); 
	% Annotations axes et unites
	\draw [axe, thick] (\xmin,0)--(\xmax,0) node[above left]  {$t$};
	\draw [axe, thick] (0,\ymin)--(0,\ymax) node[below right] {$y$};
	\foreach \x in {-3,-2,-1,0,1,2,3} 
		\draw (\x,1pt)--(\x,-1pt) node [below right]{$\x$};	
	\foreach \y in {-1,1} 
		\draw (1pt,\y)--(-1pt,\y) node [left]{$\y$};	
	
	\draw[samples=400][ color=BrickRed] [domain=-3.5:3.55][line width=1.5pt] plot[red](\x,{sin(3.14*(\x+1.3) r)});

	

	% Clip pour que les figures ne sortent pas du cadre
	%\clip (\xmin,\ymin) rectangle (\xmax,\ymax); 
\end{tikzpicture}

\acocher~~~~ fonction paire \\
\acocher~~~~ fonction impaire \\ 
\acocher~~~~ fonction périodique ~~~~ $T= ..............$ \\ 
\acocher~~~~ Aucune réponse possible \\


\begin{tikzpicture}
	% Dimensions du repere
	\def\xmin{-3.5} \def\xmax{3.55} \def\ymin{-1.25} \def\ymax{1.2}
	% Styles des axes et de la grille
	
	\tikzstyle{axe}=[->,>=stealth']
	\tikzstyle{grille}= [step=1, very thin]
	% Grille
	\draw [grille] (\xmin,\ymin) grid (\xmax,\ymax); 
	% Annotations axes et unites
	\draw [axe, thick] (\xmin,0)--(\xmax,0) node[above left]  {$t$};
	\draw [axe, thick] (0,\ymin)--(0,\ymax) node[below right] {$y$};
	\foreach \x in {-3,-2,-1,0,1,2,3} 
		\draw (\x,1pt)--(\x,-1pt) node [below right]{$\x$};	
	\foreach \y in {-1,1} 
		\draw (1pt,\y)--(-1pt,\y) node [left]{$\y$};	
	
	\draw[samples=400][ color=gray] [domain=-3.5:3.55][line width=1.5pt] plot[red](\x,{0.5*(cos((1.5*pi)*(\x) r)-0.5)});

	

	% Clip pour que les figures ne sortent pas du cadre
	%\clip (\xmin,\ymin) rectangle (\xmax,\ymax); 
\end{tikzpicture}


\acocher~~~~ fonction paire \\
\acocher~~~~ fonction impaire \\ 
\acocher~~~~ fonction périodique ~~~~ $T= ..............$ \\ 
\acocher~~~~ Aucune réponse possible \\


\end{multicols}

\begin{multicols}{2}
\begin{tikzpicture}
	% Dimensions du repere
	\def\xmin{-3.5} \def\xmax{3.55} \def\ymin{-1.25} \def\ymax{1.2}
	% Styles des axes et de la grille
	
	\tikzstyle{axe}=[->,>=stealth']
	\tikzstyle{grille}= [step=1, very thin]
	% Grille
	\draw [grille] (\xmin,\ymin) grid (\xmax,\ymax); 
	% Annotations axes et unites
	\draw [axe, thick] (\xmin,0)--(\xmax,0) node[above left]  {$t$};
	\draw [axe, thick] (0,\ymin)--(0,\ymax) node[below right] {$y$};
	\foreach \x in {-3,-2,-1,0,1,2,3} 
		\draw (\x,1pt)--(\x,-1pt) node [below right]{$\x$};	
	\foreach \y in {-1,1} 
		\draw (1pt,\y)--(-1pt,\y) node [left]{$\y$};	
	
	\draw[samples=400][ color=violet] [domain=-3.5:3.55][line width=1.5pt] plot(\x,{cos(3.14*(\x) r)*(\x*0.32)});

	

	% Clip pour que les figures ne sortent pas du cadre
	\clip (\xmin,\ymin) rectangle (\xmax,\ymax); 
\end{tikzpicture}

\acocher~~~~ fonction paire \\
\acocher~~~~ fonction impaire \\ 
\acocher~~~~ fonction périodique ~~~~ $T= ..............$ \\ 
\acocher~~~~ Aucune réponse possible \\


\begin{tikzpicture}
	% Dimensions du repere
	\def\xmin{-3.5} \def\xmax{3.55} \def\ymin{-1.25} \def\ymax{1.2}
	% Styles des axes et de la grille
	
	\tikzstyle{axe}=[->,>=stealth']
	\tikzstyle{grille}= [step=1, very thin]
	% Grille
	\draw [grille] (\xmin,\ymin) grid (\xmax,\ymax); 
	% Annotations axes et unites
	\draw [axe, thick] (\xmin,0)--(\xmax,0) node[above left]  {$t$};
	\draw [axe, thick] (0,\ymin)--(0,\ymax) node[below right] {$y$};
	\foreach \x in {-3,-2,-1,0,1,2,3} 
		\draw (\x,1pt)--(\x,-1pt) node [below right]{$\x$};	
	\foreach \y in {-1,1} 
		\draw (1pt,\y)--(-1pt,\y) node [left]{$\y$};	
	
	\draw[samples=400][ color=orange] [domain=-3.5:3.16][line width=1.5pt] plot(\x,{sin(3.14*(\x) r)*(exp(\x)*0.1)});

	

	% Clip pour que les figures ne sortent pas du cadre
	\clip (\xmin,\ymin) rectangle (\xmax,\ymax); 
\end{tikzpicture}


\acocher~~~~ fonction paire \\
\acocher~~~~ fonction impaire \\ 
\acocher~~~~ fonction périodique ~~~~ $T= ..............$ \\ 
\acocher~~~~ Aucune réponse possible \\


\end{multicols}

\newpage

\fcolorbox{black}{white}{\textbf {Exercice 2 : Savoir étudier la parité des fonctions}}\\[5 pt]
Étudier la parité des fonctions suivantes : 
\begin{multicols}{2}
\center{$f(x)=3x^2-1$ \\
\cadreblanc {40}
$g(x)=x^5+x^3+x$\\
\cadreblanc {40}
}
\end{multicols} 
\begin{multicols}{2}
\center{$h(x)=\dfrac{-5x}{2x^2+4}$ \\
\cadreblanc {37}
$j(x)=3x^2+\dfrac{1}{x^2}$\\
\cadreblanc {40}
}
\end{multicols}
\begin{multicols}{2}
\center{$k(x)=x^4-2x^2+5$ \\
\cadreblanc {40}
$l(x)=\dfrac{x^3}{x^2+1}$\\
\cadreblanc {40}
}
\end{multicols}


\end{document}